\section{Evaluation}
\label{sec:evaluation}
We ask three questions: (1) what storage does the evaluated representation
actually own initially, finally and during decoding; (2) what captioning
quality accompanies that operating point; and (3) what completed-request
time is observed after accounting for construction and reconstruction?
These questions are deliberately narrower than universal superiority at
a common byte budget or independent confirmation of the original design.

\subsection{Experimental setup and controls}
\label{sec:setup}
\paragraph{Model and hardware.}
All methods use \smallpath{llava-hf/llava-1.5-7b-hf}, model/processor revision
\texttt{b234b804}, BF16 model weights and HF eager attention
\citep{liu2024improved}. Workers use one NVIDIA A100-SXM4-80GB and 12 CPU
cores. At most six GPU workers execute concurrently; this is campaign
parallelism, not model parallelism or a batched inference claim. Other
users may occupy other GPUs on the node. Software is PyTorch 2.8.0 with
CUDA runtime 12.8, Transformers 4.55.2, lmms-eval 0.7.2 and
pycocoevalcap 1.2. Telemetry records the driver, GPU UUID, clocks, power
state and visible process concurrency; no locked-clock experiment is
claimed. The appendix lists complete revisions and settings.

\paragraph{Dataset and generation.}
We use 256 paired COCO validation images \citep{lin2014coco}, pinned through
\smallpath{lmms-lab/COCO-Caption}, revision \texttt{c3213b54}.
The prompt is ``Provide a one-sentence caption for the provided image.''
Every method receives the same image, processor/chat template,
tokenization, unpadded prompt and generation schedule. The actual decoder
prefix is 603 positions. Greedy generation computes all 128 predictions,
regardless of EOS; scoring uses the prefix ending at the first terminal
token. This separates fixed systems work from caption parsing. A final
prediction is not an additional processed/cache-appended token.

The ledger preserves every supplied reference caption, including images
with five, six or seven references. No references are truncated to force a
uniform count. The cohort was reused from the frozen storage study and
has been inspected. Historical-ID exclusions are recorded, but universal
underlying-source freshness is not certified. These are development
operating-point measurements, not untouched confirmation.

\paragraph{Methods and budgets.}
We evaluate FullKV, \Evi{}, Prompt Top-K and local PrefixKV, VL-Cache and
LOOK-M precision/storage adaptations. R+A denotes reimplementation plus
adaptation; A denotes a local adaptation. All compressed methods retain
their frozen policies. Their encoded bytes differ, despite a historical
common source-slot convention. We therefore report them as separate
operating points, not a same-budget test, swept frontier or proof that one
selection policy is intrinsically better. FullKV is an uncompressed
quality reference and is excluded from compressed-method latency tests.

\paragraph{Three separate measurement passes.}
First, 96 memory workers measure all six methods on all 256 images.
Second, 16 workers perform uninstrumented timing for those same
method--image pairs, after five warmups per method per worker and with
randomized method order within each image. Third, six fixed image indices
receive five warmups and 20 paired timed repetitions per compressed
method. Each timed request recomputes media processing, prefill,
construction, graph setup and decode; only loaded model weights and
ordinary file/compiler caching are reused. Model loading, artifact I/O and
metric scoring are excluded. Memory-audit timings are never substituted
for latency.

We use synchronized monotonic request wall time. Device events and
diagnostic intervals retain their own meanings; nested device or host
intervals are not summed as exclusive kernel costs
\citep{pytorch2025cuda}. The timer includes final CPU logits export,
EOS-aware detokenization and packet retirement. Since outputs are not
streamed, we do not claim measured time-to-first-token emission.

\subsection{Memory boundaries and correctness}
\label{sec:memory}
Initial and final storage count unique backing allocations rather than
tensor-view sizes. The decode audit additionally observes live
reconstruction/lease or native-append overlap. Its scope excludes weights,
ordinary attention activations, vendor-internal workspace and prefill
transients. Different native/encoded observers are recorded explicitly;
this is a declared decode-owned tensor boundary, not complete process HBM.

Before expansion, six canaries compare all 128 exported logits bitwise
with and without memory instrumentation. All pass, as do the measured
storage endpoints. The canonical CPU suite passes 4,548 tests. The full
study has 1,536 audited memory rows, 1,536 corresponding uninstrumented
timing rows and 600 repeated timing rows. Every audited/unaudited pair
has identical token IDs and storage endpoints. All 124 GPU workers
complete, with no failed accepted sample, missing pair or FullKV fallback.
These checks validate measurement consistency, not equivalence of a lossy
cache to unrestricted FullKV.

\begin{table*}[t]
\centering\small
\setlength{\tabcolsep}{6pt}
% Exported from hash-verified measured records; do not hand edit.
\begin{tabular}{lrrrrrr}
\toprule
Method & Initial & \shortstack{Decode-owned peak\\mean/max} & Final & CIDEr $\uparrow$ & p50 $\downarrow$ & p95 $\downarrow$\\
\midrule
FullKV reference & 301.50 & 370.70/370.70 & 365.00 & 116.83 & 5.85 & 5.97\\
\textbf{EviCastKV-S} & 41.14 & 83.52/84.24 & 73.22 & 119.77 & 6.55 & 6.72\\
Prompt Top-K (A) & 50.73 & 92.83/92.95 & 82.80 & 118.33 & 7.50 & 7.66\\
PrefixKV (R+A) & 45.70 & 86.01/86.56 & 77.78 & 103.10 & 7.38 & 7.54\\
VL-Cache (R+A) & 53.31 & 94.39/94.64 & 85.39 & 122.16 & 13.25 & 17.66\\
LOOK-M (R+A) & 50.72 & 92.82/92.83 & 82.80 & 120.96 & 48.86 & 51.59\\
\bottomrule
\end{tabular}

\caption{Completed 256-image COCO operating points: initial/final memory
and peak mean/max in MiB; full-request latency quantiles in seconds. Memory
columns are means except the explicitly shown maximum. R+A means local
reimplementation plus precision/storage adaptation, not official execution.
\textbf{Realized byte budgets differ.} FullKV latency is descriptive only.
No row is omitted for unfavorable quality or latency.}
\label{tab:main}
\end{table*}

\paragraph{Observed storage.}
Table~\ref{tab:main} reports the complete cohort. \Evi{} owns
\EviInitial~MiB initially, \EviFinal~MiB finally and
\EviPeak~MiB mean audited decode-owned peak, with a maximum of
\EviMaxPeak~MiB. Relative to FullKV, those mean boundaries decrease by
\FullInitialReduction\%, \FullFinalReduction\% and \FullPeakReduction\%,
respectively. These KV-only reductions must not be applied to total GPU
memory, which also contains the model and other resources.

Relative to Prompt Top-K's observed point, mean initial/final/peak
reductions are \TopInitialReduction\%, \TopFinalReduction\% and
\TopPeakReduction\%. PrefixKV's adaptation is the closest measured
compressed memory point: its final and mean decode-peak differences are
only 5.85\% and 2.90\%. Generated history therefore narrows the initial
memory advantage. No comparison here establishes a 30\% reduction against
the compressed competitors, much less a 30\% equal-quality frontier shift.

\subsection{Corpus CIDEr and paired uncertainty}
\label{sec:quality}
CIDEr measures reference-caption consensus \citep{vedantam2015cider}.
We use the pinned PTB tokenizer and corpus reducer through the COCO
evaluation ecosystem, report score points on the evaluator's $\times100$
scale, and retain all images in the denominator. No mean of shard CIDEr
scores is used. The bootstrap samples paired images with replacement
10,000 times using seed 20261106. Each replicate recomputes reference
document frequencies, TF--IDF weights and the complete corpus reduction
\citep{efron1979bootstrap}. Repeated images count as repeated resampled
documents; cached per-image scores are not substituted.

\begin{table*}[t]
\centering\small
% Exported from hash-verified measured records; do not hand edit.
\begin{tabular}{lrrr}
\toprule
Reference point & $\Delta$CIDEr & Paired 95\% interval & Familywise 95\% interval\\
\midrule
FullKV reference & +2.94 & [+0.005, +6.124] & [-0.915, +7.078]\\
Prompt Top-K (A) & +1.44 & [-1.726, +4.749] & [-2.713, +5.712]\\
PrefixKV (R+A) & +16.67 & [+7.429, +24.811] & [+4.677, +27.616]\\
VL-Cache (R+A) & -2.38 & [-6.386, +1.424] & [-7.805, +2.533]\\
LOOK-M (R+A) & -1.18 & [-5.407, +2.913] & [-6.670, +4.182]\\
\bottomrule
\end{tabular}

\caption{\Evi{} minus each reference operating point in corpus CIDEr.
Five contrasts form the quality family. Unadjusted paired intervals and
Bonferroni familywise 95\% intervals are both shown. These intervals quantify
paired uncertainty on the inspected cohort; they do not establish dataset
freshness or equal-byte comparisons.}
\label{tab:quality}
\end{table*}

\Evi{} scores \EviCider{} versus \FullCider{} for FullKV and 118.33 for
Prompt Top-K. Its Top-K difference is $+1.44$ CIDEr, but the adjusted
interval includes negative values. The VL-Cache and LOOK-M adaptations
score 122.16 and 120.96, respectively, both above \Evi{}'s point estimate.
Table~\ref{tab:quality} therefore does not establish zero-margin quality
non-inferiority to those compressed points. An interval crossing zero is
not evidence of equality.

Only the PrefixKV-adaptation difference remains positive after this
familywise correction. That finding requires an additional caveat:
42/256 of its outputs reach the fixed 128-prediction limit, versus none
for the other five methods. All remain included and scored; none are
silently dropped. This behavior, the added codec and the changed task mean
that the result cannot be interpreted as a reproduction failure or general
quality claim about official PrefixKV. Appendix tables report caption
lengths and limit-hit counts explicitly.

\begin{figure}[t]
\centering\begin{tikzpicture}
\begin{axis}[
 width=\columnwidth,height=5.6cm,
 xlabel={Initial encoded storage (MiB)},ylabel={Corpus CIDEr},
 xmin=39,xmax=56,ymin=99,ymax=125,
 grid=major,grid style={gray!15},tick label style={font=\scriptsize},
 label style={font=\small},
 legend style={at={(.5,1.04)},anchor=south,legend columns=2,draw=none,font=\scriptsize},
 legend cell align=left]
\addplot[only marks,mark=star,mark size=3pt,color=eviblue]
 table[x=initial,y=cider]{generated/point_EviCastKV.dat};
\addlegendentry{EviCastKV-S}
\addplot[only marks,mark=square*,mark size=2.4pt,color=exactgreen]
 table[x=initial,y=cider]{generated/point_TopK.dat};
\addlegendentry{Prompt Top-K (A)}
\addplot[only marks,mark=triangle*,mark size=2.6pt,color=patchorange]
 table[x=initial,y=cider]{generated/point_PrefixKV.dat};
\addlegendentry{PrefixKV (R+A)}
\addplot[only marks,mark=diamond*,mark size=2.5pt,color=riskred]
 table[x=initial,y=cider]{generated/point_VLCache.dat};
\addlegendentry{VL-Cache (R+A)}
\addplot[only marks,mark=*,mark size=2.3pt,color=violet]
 table[x=initial,y=cider]{generated/point_LOOKM.dat};
\addlegendentry{LOOK-M (R+A)}
\end{axis}
\end{tikzpicture}

\caption{Discrete compressed operating points, not curves or a fitted Pareto
frontier. \Evi{} has the smallest initial cache here, but VL-Cache and LOOK-M
adaptations have higher point-estimate CIDEr. FullKV is outside the plotted
compressed-memory range and is listed in Table~\ref{tab:main}.}
\label{fig:quality}
\end{figure}

\subsection{Full-request time and observed execution policy}
\label{sec:latency}
On the 256 images, \Evi{}'s median/p95 request time is
\EviMedian/\EviTail~s. This is lower than the measured compressed adaptations,
but these are not matched-byte or official-code latency conclusions.
In the separate six-image repeated study, ratios are recomputed inside a
nested paired bootstrap over images and repetitions. Four contrasts use
familywise-adjusted intervals. The ratio of mean \Evi{} time to mean
baseline time is below one for all four observed points, yielding
\TopFaster\%, \PrefixFaster\%, \VLFaster\% and \LookFaster\% lower time
versus Top-K, PrefixKV, VL-Cache and LOOK-M adaptations, respectively.
The six-image ratios are not pooled with the 256-image timing distribution.

\begin{figure}[t]
\centering\begin{tikzpicture}
\begin{axis}[
 width=.94\columnwidth,height=4.6cm,
 xlabel={EviCastKV-S / baseline request time},xmin=0,xmax=1.06,
 symbolic y coords={LOOKM,PrefixKV,TopK,VLCache},
 ytick=data,yticklabels={LOOK-M (R+A),PrefixKV (R+A),Prompt Top-K (A),VL-Cache (R+A)},
 enlarge y limits=.25,grid=major,grid style={gray!15},
 tick label style={font=\scriptsize},label style={font=\small}]
\addplot[only marks,color=eviblue,mark=*,mark size=2.5pt,
 error bars/.cd,x dir=both,x explicit]
 table[x=ratio,y=baseline,x error minus=minus,x error plus=plus]{generated/latency.dat};
\draw[dashed,color=riskred] (axis cs:1,LOOKM) -- (axis cs:1,VLCache);
\end{axis}
\end{tikzpicture}

\caption{Ratio of mean completed-request time on the six-image, 20-repeat
study. Bars show multiplicity-adjusted paired intervals; the reference line
is equal time. FullKV is excluded. Different realized cache layouts and
different eligible execution paths remain part of these operating points.}
\label{fig:latency}
\end{figure}

\paragraph{What explains the timing, and what remains unproven?}
Saved phase boundaries show approximately 0.80~s for \Evi{}'s combined
prefill/construction/retirement and 5.70~s for its recorded decode region.
Top-K has a similar 0.80~s first region and a 6.64~s decode region. This is
consistent with the importance of reconstruction scheduling and replay,
not proof that fewer bytes alone cause the gain: \Evi{} engages uniform
replay on every image while the frozen baseline layouts are ineligible.
VL-Cache's local adaptation spends approximately 5.34~s in its combined
prefill/construction region; LOOK-M's spends 42.05~s. Those are properties
of these local paths, not intrinsic costs of the published algorithms.
The appendix gives complete phase means and engagement counts.

The broad six-component planner is inactive, so its utility, graph and
forecasting cannot explain these numbers. Nor can shorter-context eager
results establish speedups under paged/FlashAttention serving engines,
long contexts or larger batches \citep{dao2022flashattention,kwon2023paged}.
